% ============================================
% DOC-02 (EN)
% English translation — SPEC-935-R355
% ============================================
\section*{\textbf{DOC-02} --- Medical Record of Patient 1,260}
\markboth{DOC-02 --- Medical Record of Patient 1,260}{}

\noindent \textit{Barbacena Hospital-Colony --- Admission Form}

\noindent \textbf{Barbacena Hospital-Colony --- Admission Form}
\textit{Fictional reconstruction of an admission form; this is not an authentic archival document.}
\textit{Narrative composition based on historical context, without endorsement by the Minas Gerais Public Archive.}

\vspace{0.5em}\noindent\rule{\textwidth}{0.5pt}\vspace{0.5em}

\noindent \textbf{Registration No.:} 1,260  
\textbf{Name:} Joaquim Antônio Correia  
\textbf{Date of Birth:} c. 1907 (estimated)  
\textbf{Place of Birth:} Quixeramobim, Ceará  
\textbf{Parentage:} Father: Antônio Correia (deceased --- starvation, 1915); Mother: Maria do Rosário Correia (missing --- forced relocation, 1915)  
\textbf{Origin:} Retirante who left Senador Pompeu in 1915 and was confined at the Alagadiço Camp, near Fortaleza. Transferred in 1917 via Santa Casa de Misericórdia of Belo Horizonte.

\vspace{0.5em}\noindent\rule{\textwidth}{0.5pt}\vspace{0.5em}

\noindent \textbf{Admission Examination --- Dr. Nelson Fontes, November 12, 1917:}

\noindent \textit{\textquotedblleft{}Patient presents in a state of severe cachexia. Estimated weight: 20 kg (expected for age: 28-30 kg). Evident protein-calorie malnutrition. Edema in the lower limbs. Distended abdomen. Compromised dentition (multiple caries, loss of 3 teeth).}

\noindent \textit{Most notable sign: heterochromia of the right iris, which presents in a golden-yellow tone, absent on the left. Patient claims that 'a beast' changed its color. Probable infectious or traumatic sequela, unidentified.}

\noindent \textit{Mental state: placid, cooperative, but with a fixed, vacant gaze. Answers questions in monosyllables. When questioned about his family, he declared: 'They died. All of them.' No signs of sensorium alteration. Oriented in time and space.}

\noindent \textit{Preliminary assessment: Hebephrenic schizophrenia (dementia praecox). Isolation and observation recommended.\textquotedblright{}}

\vspace{0.5em}\noindent\rule{\textwidth}{0.5pt}\vspace{0.5em}

\noindent \textbf{Incident History:}

\begin{moltablethree}
\hline
Date & Incident & Consequence \\
\hline
Feb/1920 & Unprovoked assault on an orderly & 30 days of isolation \\
Jun/1922 & \textbf{Homicide --- Orderly Wagner} & 90 days of isolation \\
Mar/1937 & \textbf{Homicide --- Patient from Ward C} & Permanent isolation \\
Apr/1937 & Self-mutilation (tore out nails with his teeth) & Physical restraint \\
Sep/1945 & Episode of uncontrollable vocalization (claims 'the Shroud-Ripper' is speaking through him) & Bromide sedation \\
Jan/1948 & \textbf{Homicide --- Visitor (entered the forbidden ward)} & Perpetual isolation \\
1949-1956 & Selective mutism (7 years without speaking to staff) & --- \\
1956 & Return of speech: first phrase 'The diary' & --- \\
1979 & Sessions with Dr. Heitor Oliveira & See attached reports \\
\hline
\end{moltablethree}

\vspace{0.5em}\noindent\rule{\textwidth}{0.5pt}\vspace{0.5em}

\noindent \textbf{Medications Administered (partial record):}

\noindent - Potassium bromide (1920-1945) --- sedative
- Electroshock (1937, 1945, 1952) --- without anesthesia
- Transorbital lobotomy --- \textit{indicated in 1948 but not performed (patient reacted with extreme violence to the mere attempt to approach)}
- Chlorpromazine (1955 onward) --- antipsychotic

\vspace{0.5em}\noindent\rule{\textwidth}{0.5pt}\vspace{0.5em}

\noindent \textbf{Final Remarks:}

\noindent Patient 1,260 is one of the Colony's oldest inmates. His behavior is unpredictable: periods of absolute calm interspersed with outbursts of extreme violence. He keeps a diary that he never allowed any staff member to read during his lifetime.

\noindent \textit{Additional note (1979): Dr. Heitor Oliveira obtained permission to read the diary. First and only person to have access to the document. Dr. Oliveira requested medical leave 48 hours later. He never returned to work.}

\noindent \textit{The diary is currently under the custody of the Barbacena coroner's institute, attached to the autopsy report.}

\vspace{0.5em}
\noindent\textit{\textbf{↪ Routes:} \rota{MEM-08} \rotasep \rota{DOC-03} \rotasep \rota{LUC-01}}
